\SourceLeaf{146}%
\Article{64}
Proficiscendo a regula ultima in art. praec. eruta invenimus esse
\[
\begin{array}{c|l}
\text{numeri}&\text{characterem}\equiv\\\hline
-1+i&\frac18(aa+2ab-bb-1)\\
-1-i&\frac18(-aa+2ab+bb+1)\\
+1-i&\frac18(aa+2ab+3bb-1)
\end{array}
\]

Hoc statim inde sequitur, quod character ipsius $i$ est $\frac14(aa+bb-1)$, character ipsius $-1$ autem $\frac12(aa+bb-1)\equiv\frac12bb$, quum $aa-1$ semper sit formae $8n$. Manifesto hae quatuor regulae, etiamsi hactenus ab inductione mutuatae sint, ita inter se sunt nexae, ut quamprimum unius demonstratio absoluta fuerit, tres reliquae simul sint demonstratae. Vix opus est monere, etiam in his regulis tantummodo supponi $a$ imparem, $b$ parem.

Si formulas ad modulos primarios restrictas adhibere non displicet, hac forma uti possumus. Est
\[
\begin{array}{c|l}
\text{numeri}&\text{character}\equiv\\\hline
-1+i&\frac14(-a-b+1-bb)\\
-1-i&\frac14(a-b-1+bb)\\
+1-i&\frac14(-a-b+1+bb)
\end{array}
\]

Formulae simplicissimae prodeunt, si, ut initio inductionis nostrae feceramus, modulos primi et secundi generis distinguimus. Est scilicet character
\[
\begin{array}{c|c|c}
\text{numeri}&\text{pro modulis primi generis}&\text{pro modulis secundi generis}\\\hline
-1+i&\frac14(-a-b+1)&\frac14(-a-b-3)\\
-1-i&\frac14(a-b-1)&\frac14(a-b+3)\\
+1-i&\frac14(-a-b+1)&\frac14(-a-b+5)
\end{array}
\]

\Article{65}
Pro numero $-1+2i$, ad quem iam progredimus, eandem distinctionem inter modulos $a+bi$ eos, pro quibus $a\equiv1,b\equiv0$, atque eos, pro quibus $a\equiv3,b\equiv2$ quoque adhibebimus. Tabula art.~62 docet, respectu illius numeri respondere
\SourceLeaf{147}%
\[
\begin{array}{c|l}
\text{characterem}&\text{modulis primi generis}\\\hline
0&-3+8i,\ +5-8i,\ +9+4i,\ -11+4i\\
1&+1+4i,\ +5-4i,\ -7,\ -3-8i\\
2&+1-4i,\ +5+8i,\ -7-8i,\ -11\\
3&-3,\ +5+4i,\ +9-4i,\ -7+8i,\ -11-4i
\end{array}
\]

Revocatis singulis his modulis ad residua absolute minima secundum modulum $-1+2i$, animadvertimus, omnes, quibus respondet character $0$, esse $\equiv1$; eos, quibus character $1$ respondet, $\equiv i$; eos, quorum character est $2$, fieri $\equiv-1$; denique omnes, quorum character est $3$, fieri $\equiv-i$. At characteres numerorum $1,i,-1,-i$ pro modulo $-1+2i$ ipsi sunt $0,1,2,3$ resp.; quapropter in omnibus his $17$ exemplis character numeri $-1+2i$ respectu moduli prioris generis $a+bi$, cum charactere huius numeri respectu moduli $-1+2i$ identicus est.

Perinde adiumento tabulae invenitur, respondere
\[
\begin{array}{c|l}
\text{characterem}&\text{modulis secundi generis}\\\hline
0&+3+2i,\ -5-2i,\ -1+10i,\ -1-10i,\ +11+6i\\
1&+3-2i,\ -1+6i,\ -5-6i,\ +7+10i,\ +7-10i\\
2&-5+2i,\ -1-6i,\ +7-2i\\
3&-1-2i,\ +7+2i,\ -5+6i,\ +3+10i,\ +3-10i,\ +11-6i
\end{array}
\]

Revocatis his modulis ad residua minima secundum modulum $-1+2i$, omnia, quibus resp. characteres $0,1,2,3$ respondent, congrua inveniuntur numeris $-1,-i,+1,+i$; his vero ipsis numeris, si vice versa $-1+2i$ pro modulo adoptatur, competent characteres $2,3,0,1$ resp. Quapropter in omnibus his $19$ exemplis character numeri $-1+2i$ respectu moduli secundi generis duabus unitatibus differt a charactere huius numeri respectu numeri $-1+2i$ pro modulo habiti.

Ceterum nullo negotio perspicitur, prorsus similia respectu numeri $-1-2i$ locum habitura esse.

\Article{66}
Pro numero $-3$ distinctionem inter modulos primi generis et secundi omittimus, quum eventus doceat, illam hic superfluam esse. Respondet itaque
\SourceLeaf{148}%
\[
\begin{array}{c|l}
\text{character}&\text{modulis}\\\hline
0&\begin{gathered}-1+6i,\ -1-6i,\ -7,\ -5+6i,\\-5-6i,\ -11,\ 11+6i,\ 11-6i\end{gathered}\\[3pt]
1&\begin{gathered}-1-2i,\ 1-4i,\ -5+2i,\ 5+4i,\ 7+2i,\\5-8i,\ -1+10i,\ -7-8i,\ -11-4i,\ 7-10i\end{gathered}\\[3pt]
2&\begin{gathered}3+2i,\ 3-2i,\ -3+8i,\ -3-8i,\\9+4i,\ 3+10i,\ 3-10i\end{gathered}\\[3pt]
3&\begin{gathered}-1+2i,\ 1+4i,\ -5-2i,\ 5-4i,\ 7-2i,\\5+8i,\ -1-10i,\ -7+8i,\ -11+4i,\ 7+10i\end{gathered}
\end{array}
\]

Revocatis his modulis ad residua minima secundum modulum $3$, videmus, eos, quibus respondet character $0$, esse partim $\equiv1$, partim $\equiv-1$; eos, quorum character est $1$, fieri vel $\equiv1-i$, vel $\equiv-1+i$, eos, quorum character est $2$, fieri vel $\equiv i$, vel $\equiv-i$; denique eos, quibus competit character $3$, esse vel $\equiv1+i$, vel $\equiv-1-i$. Ex hac itaque inductione colligimus, characterem numeri $-3$ pro modulo, qui est numerus primus inter associatos primarius, identicum esse cum charactere huius ipsius numeri, dum $3$, sive, quod eodem redit, $-3$ tamquam modulus consideratur.

\Article{67}
Simili inductione circa alios numeros primos instituta, invenimus, numeros $3\pm2i,-1\pm6i,7\pm2i,-5\pm6i$ etc. suppeditare theoremata ei similia, ad quod in art.~65 respectu numeri $-1+2i$ pervenimus; contra numeros $1\pm4i,5\pm4i,-3\pm8i,5\pm8i,9\pm4i$ etc. perinde se habere ut numerum $-3$. Inductio itaque perducit ad elegantissimum theorema,. quod ad instar theoriae residuorum quadraticorum in arithmetica numerorum realium \textsc{Theorema fundamentale} theoriae residuorum biquadraticorum nuncupare liceat, scilicet:

\emph{Denotantibus $a+bi,a'+b'i$ numeros primos diversos inter associatos suos primarios, i.~e. secundum modulum $2+2i$ unitati congruos, character biquadraticus numeri $a+bi$ respectu moduli $a'+b'i$ identicus erit cum charactere numeri $a'+b'i$ respectu moduli $a+bi$, si vel uterque numerorum $a+bi,a'+b'i$, vel alteruter saltem, ad primum genus refertur, i.~e. secundum modulum $4$ unitati congruus est: contra characteres illi duabus unitatibus inter se different, si neuter numerorum $a+bi,a'+b'i$ ad primum genus refertur, i.~e. si uterque secundum modulum $4$ congruus est numero $3+2i$.}

\SourceLeaf{149}%
At non obstante summa huius theorematis simplicitate, ipsius demonstratio inter mysteria arithmeticae sublimioris maxime recondita referenda est, ita ut, saltem ut nunc res est, per subtilissimas tantummodo investigationes enodari possit, quae limites praesentis commentationis longe transgrederentur. Quamobrem promulgationem huius demonstrationis, nec non evolutionem nexus inter hoc theorema atque ea, quae in initio huius commentationis per inductionem stabilire coeperamus, ad commentationem tertiam nobis reservamus. Coronidis tamen loco iam hic trademus, quae ad demonstrationem theorematum in art.~63.~64 propositorum requiruntur.

\Article{68}
Initium facimus a numeris primis $a+bi$ talibus, pro quibus $b=0$ (tertia specie art.~34), ubi itaque (ut numerus inter associatos primarius sit) $a$ debet esse numerus primus realis negativus formae $-(4n+3)$, pro quo scribemus $-q$, quales sunt $-3,-7,-11,-19$ etc. Denotando per $\lambda$ characterem numeri $1+i$, illo numero pro modulo accepto, esse debet
\[
i^\lambda\equiv(1+i)^{\frac14(qq-1)}\equiv2^{\frac18(qq-1)}\pdot i^{\frac18(qq-1)}\pmod q
\]
Sed constat, $2$ esse residuum quadraticum, vel non-residuum quadraticum ipsius $q$, prout $q$ sit formae $8n+7$, vel formae $8n+3$, unde colligimus, esse generaliter
\[
2^{\frac12(q-1)}\equiv(-1)^{\frac14(q+1)}\equiv i^{\frac12(q+1)}\pmod q
\]
adeoque evehendo ad potestatem exponentis $\frac14(q+1)$
\[
2^{\frac18(qq-1)}\equiv i^{\frac18(q+1)^2}\pmod q
\]
Aequatio itaque praecedens hanc formam induit
\[
i^\lambda\equiv i^{\frac18(q+1)^2+\frac18(qq-1)}\equiv i^{\frac14(qq+q)}\pmod q
\]
unde sequitur
\[
\lambda\equiv\frac14(qq+q)\equiv\frac14(q+1)^2-\frac14(q+1)\pmod4
\]
sive quum habeatur $\frac14(q+1)^2\equiv0\pmod4$, $\lambda\equiv-\frac14(q+1)\equiv\frac14(a-1)\pmod4$. Quod est ipsum theorema art.~63 pro casu $b=0$.

\SourceLeaf{150}%
\Article{69}
Longe vero difficilius absolvuntur moduli $a+bi$ tales, pro quibus non est $b=0$ (numeri quartae speciei art.~34), pluresque disquisitiones erunt praemittendae. Normam $aa+bb$, quae erit numerus primus realis formae $4n+1$, designabimus per $p$.

Denotetur per $S$ complexus omnium residuorum simpliciter minimorum pro modulo $a+bi=m$, exclusa cifra, ita ut multitudo numerorum in $S$ contentorum sit $=p-1$. Designet $x+yi$ indefinite numerum huius systematis, statuaturque $ax+by=\xi,ay-bx=\eta$. Erunt itaque $\xi,\eta$ integri inter limites $0$ et $p$ \emph{exclusive} contenti; in casu praesente enim, ubi $a,b$ inter se primi sunt, formulae art.~45, puta $\eta\equiv k\xi,\xi\equiv-k\eta\pmod p$ docent, neutrum numerorum $\xi,\eta$ esse posse $=0$, nisi alter simul evanescat, adeoque fiat $x=0,y=0$, quam combinationem iam eiecimus. Criterium itaque numeri $x+yi$ in $S$ contenti, consistit in eo, ut quatuor numeri $\xi,\eta,p-\xi,p-\eta$ sint positivi.

Praeterea observamus pro nullo tali numero esse posse $\xi=\eta$; hinc enim sequeretur $p(x+y)=a(\xi+\eta)+b(\xi-\eta)=2a\xi$, quod est absurdum, quum nullus factorum $2,a,\xi$ per $p$ divisibilis sit. Simili ratione aequatio $p(x-y+a+b)=2a\xi+(a+b)(p-\xi-\eta)$ docet, esse non posse $\xi+\eta=p$. Quapropter quum numeri $\xi-\eta,p-\xi-\eta$ esse debeant vel positivi vel negativi, hinc petimus subdivisionem systematis $S$ in quatuor complexus $C,C',C'',C'''$, puta ut coniiciantur
\[
\begin{array}{c|ll}
\text{in complexum}&\multicolumn{2}{c}{\text{numeri pro quibus}}\\\hline
C&\xi-\eta\text{ positivus},&p-\xi-\eta\text{ positivus}\\
C'&\xi-\eta\text{ positivus},&p-\xi-\eta\text{ negativus}\\
C''&\xi-\eta\text{ negativus},&p-\xi-\eta\text{ negativus}\\
C'''&\xi-\eta\text{ negativus},&p-\xi-\eta\text{ positivus}
\end{array}
\]

Criterium itaque numeri complexus $C$ proprie sextuplex est, puta sex numeri $\xi,\eta,p-\xi,p-\eta,\xi-\eta,p-\xi-\eta$ positivi esse debent; sed manifesto conditiones $2,5$ et $6$ iam sponte implicant reliquas. Similia circa complexus $C',C'',C'''$ valent, ita ut criteria completa sint triplicia, puta
